\chapter{Reading the experimental sequence}\label{ch:evidence-route}
We now know enough about the historical engine to read its experiments
without confusing their object with the current Gaussian proposal. The
sequence matters. Early tests checked local mechanisms and numerical
behavior. Later work exposed limitations of the questions being asked.
Quote conformance and service-sensitive comparisons then required different
records and different interpretations. A list of green tests would erase
that development.

\section{Begin with the question, not the result}
For every experimental chapter, write down the question in one sentence
before reading its outcome. Is the test about conservation, a numerical
certificate, exact quotation, execution performance, preservation, service
or the ability to discriminate alternatives? These are not interchangeable
targets. A fast evaluator can be scientifically unhelpful; a preservation
policy can correctly refuse useful service; an exact quote can be embedded
in a controller that never chooses the action being quoted.

A result is meaningful when its design could have exposed the failure
relevant to its question. A stress case that almost never reaches a hard
boundary says little about behavior at that boundary. Conversely, a case
designed to expose a boundary failure should not be discarded as unrealistic
merely because it finds one. The registered scope governs the inference.

\section{Mechanism tests and dynamical claims}
The D-series belongs to a specific historical plant and response mechanism.
Its first chapters teach the value of examining conservation, loss handling,
directional flow and numerical behavior separately. An invariant can hold
while another property fails. This is not an inconsistency: the properties
have different assumptions and different failure modes.

For example, a synchronous implementation can preserve mass exactly and
still overshoot the minimum of a convex potential in one large step. A
soft penalty can strongly discourage a boundary crossing without making it
impossible. A hard source certificate can forbid an export without solving
the destination's unmet need. These logical distinctions remain important
even when the active potential family changes.

The original D9 and D10 discussions are therefore retained in full. Their
purpose here is not to reintroduce old thresholds into the Gaussian state.
It is to make the difference between a cost, a physical condition and an
admission rule concrete. Removing their vocabulary without preserving the
lesson would allow the same category error to return under new symbols.

\section{Performance is not efficacy}
Historical D5 includes software performance work. A timing measurement can
support a claim about that workload on that environment. It cannot by itself
show that the selected physical actions maintain function, nor can it price
an unregistered future cloud campaign. The input sizes, host conditions and
measured operation are part of the performance claim.

This boundary is not a reason to omit performance work. A scientifically
meaningful method may still be impractical to execute or audit. The right
conclusion is to report feasibility and scientific effect as separate
dimensions, not to allow one to stand in for the other.

\section{Quotation conformance asks a narrower question}
Gate 1B concerns the relationship between an exact local quote and its
declared mathematical and implementation contract. Such work can detect a
missing factor, an inconsistent baseline or an incorrect finite schedule.
It does not establish that a controller will choose useful actions over a
long horizon. Its success is necessary for trusting later value records,
not sufficient for believing an institutional story about them.

In the Gaussian revision, some formulae become simpler because curvature
is constant. The need to bind a quote to an epoch, support, parameter
version and accepted quantity does not disappear. Historical conformance
remains evidence about the historical implementation; new conformance must
be demonstrated for a separately accepted implementation.

\section{When a question must change}
Adversarial search can expose a defect or a limitation in the claim being
tested. A subsequent study may legitimately ask a better question. The
scientific obligation is to record that change prospectively, not pretend
that the revised interpretation was always the original criterion.

Service alignment is an especially important example because preservation
alone can favor doing nothing. An idle controller may keep every stock
unchanged while providing no requested service. Whether that is success
depends on the registered question, not on the fact that one physical
invariant remains perfect. The new first Gaussian domain has no service
objective hidden inside its random selector; service claims would require
their own represented quantities and comparison.

\section{A worksheet for every historical result}
Read each result through six questions. What was fixed before execution?
What implementation actually ran? What observation was recorded? Which
alternative interpretation was ruled out? Which limitation remains? What
later claim would exceed the evidence?

This worksheet distinguishes three situations that can otherwise look
similar. A correctly completed null result can mean that the arms do not
differ under the registered metric. An inconclusive result can mean the
design cannot decide the claim. An invalid result can mean the execution
did not satisfy its contract. None should be relabelled as another merely
to make the research narrative smoother.

The following source reading preserves the experimental progression and
its original qualifications. The next chapter then treats the original
flat O14 outcome and the later finalized Gate 1D-C evidence as separate
records. Neither is a substitute for an unperformed Gaussian experiment.
